

Unfortunately, the expected loss of a strategy is a functional of the loss function.
Therefore, based on the principle of objectivity, we should not directly use the expected loss of the optimal Bayesian
strategy for evaluating probability measures.
We need a measure that quantifies the similarity between two probability measures while it ignores the
specific loss function we are using.

Suppose that for an environment, we have decided on a probabilistic model \M, and without any particular prior
knowledge, we have chosen a probability measure $P$ from our model.
In particular, we are interested in the distribution of a random variable \X.

Imagine that we ask an expert---a person familiar with the environment---to evaluate our probability measure $P$.
Presumably, that expert has his own beliefs about the environment, and thinks that we should, instead of $P$, use a
different probability measure $Q$.
When evaluating $P$, he would try to measure the loss of using the incorrect distribution $P(\X)$,
while he believes the correct distribution is $Q(\X)$.

Let us assume that a real number can be used to measure this loss, and denote it as $\D[\X]$.
It seems reasonable, as $P$ gets closer to $Q$ this measure decreases, and increases when $P$ diverges from $Q$.
Therefore, $\D$ effectively is a measure of the divergence of $P(\X)$ from $Q(\X)$, when $Q$ is
the believed true probability measure.

Interestingly, if we make certain assumptions about the properties of this divergence measure,
it turns out that---up to a constant factor---it must be equal to the Kullback–Leibler (KL) divergence.

Alternatively, by adopting a more minimal set of properties, we can prove that if we accept Shannon’s entropy as a
measure of the uncertainty of a distribution, we should also consider KL divergence as the appropriate method for
measuring divergence between distributions.

\begin{theorem}
    Let \X be a discrete random variable.
    Suppose that $\D[\X]$ is a divergence measure that quantifies the divergence of $P(\X)$ from $Q(\X)$.
    Consider the following properties:
    \begin{description}
        \item[Consistency with entropy] Let $\varOmega$ be a continuous random variable, and $Q$ be a probability
        measure.
        If $U(\varOmega)$ is the uniform distribution with the same support as $Q(\varOmega)$, then we have
        \[
            \D[\varOmega][Q][U] = \entropy{U}{\varOmega} - \entropy{Q}{\varOmega},
        \]
        where $\entropy{}{}$ denotes Shannon's entropy.
        \item[Additivity] For any two discrete random variables $\XY$, and probability measures $P,Q$
        \[
            \D[\XY] = \D + \sumx \D[\YIx]Q(\rval[]{x})\cdot
        \]
    \end{description}
    If \measure{D} satisfies these properties, then it is the Kullback-Leibler divergence.
\end{theorem}

\begin{proof}
    \newcommand{\T}{\varOmega}
    \renewcommand{\t}{\omega}

    \newcommand{\TIx}{\T \mid \rval[]{x}}
    \newcommand{\TX}{\T, \X}
    \newcommand{\tx}{\t, \rval[]{x}}
    \newcommand{\tent}[1]{\int q(#1) \ln q(#1) d\t}

    \newcommand{\Qx}{Q(\rval[]{x})}
    Let $P(\X)$ and $Q(\X)$ be two arbitrary distributions for a discrete random variable \X.
    We define a new continuous random variable $\T$ and let
    \begin{equation*}
        p(\T = \t \mid \rvEq) =
        \begin{cases}
            \frac{1}{P(\rval[]{x})} & 0 \leq \t \leq P(\rvEq), \\
            0                       & \text{otherwise}\cdot
        \end{cases}
    \end{equation*}
    Notice that both $P(\T \mid \X)$ and $P(\TX)$ are uniform distributions, and
    $\entropy{P}{\TIx} = \ln P(\rval[]{x})$, $\entropy{P}{\TX} = 0$.

    We define $Q(\TIx)$ to be an arbitrary distribution that has the same support as $P(\TIx)$.
    This ensures that $Q(\TX)$ and $P(\TX)$ have the same support as well.
    Thus, we have
    \begin{align*}
        \D[\TX] & = \entropy{P}{\TX} - \entropy{Q}{\TX}    \\
        & = \sumx \tent{\tx}                       \\
        & = \sumx \Qx \int q(\tIx) \ln q(\tx) d\t,
    \end{align*}
    and
    \begin{align*}
        \D[\TIx] & = \entropy{P}{\TIx} - \entropy{Q}{\TIx} \\
        & = \ln P(\rval[]{x}) + \tent{\tIx}\cdot
    \end{align*}
    Now, we use the additivity property to obtain a formula for our divergence measure:
    \begin{align*}
        \D & = \D[\TX] - \sumx \D[\TIx]\Qx                                                              \\
        & = \sumx \Qx  \left( \int q(\tIx) \ln \frac{q(\tx)}{q(\tIx)}d\t - \ln P(\rval[]{x}) \right) \\
        & = \sumx \Qx  \left( \ln \Qx \int q(\tIx) d\t - \ln P(\rval[]{x}) \right)                   \\
        & = \sumx \Qx \ln \frac{\Qx}{P(\rval[]{x})}\cdot \qedhere
    \end{align*}
\end{proof}

In this theorem, the first property simply states that the entropy of a distribution should be related to the
divergence of the uniform distribution from it.
The more the divergence is, the less the entropy must be.

The second property tells us that the divergence is an additive quantity.
We know that the divergence of $P(\XY)$ from $Q(\XY)$ is a measure of the generic loss of using $P(\XY)$
when we believe $Q(\XY)$ is the true distribution.
To calculate this loss, we can assume that \X happens before \Y.
In that way, this loss would include the loss of using $P(\X)$, and then, if we observe
$\assign{X}$, which can happen with the probability $Q(\x)$, it would also include the loss of
using $P(\Y \mid \assign{x})$.
This property indicates that for combining these individual terms we should use addition.

Thus, if $Q$ is the probability measure that represents our beliefs, we can use $\D$ to evaluate any other
probability measure $P$, with respect to a random variable \X.
In other words, we evaluate $P(\X)$ by $Q(\X)$.

If we have two random variables, \X and \Y, there would be two distinct cases that should not be
confused:
We may want to evaluate $P(\XY)$, or we may want to evaluate $P(\X)$ and $P(\Y)$ separately.
Evaluating the joint distribution $P(\XY)$ is straightforward and can be done by defining a new random variable
$\rv{Z} = (\XY)$.

The second case, however, is more challenging and the correct approach is not immediately clear.
Here, we may use a similar argument to what we used for the additive property:
Suppose that we are evaluating the loss of using $P(\XY)$, when \X and \Y are independent.
If \X happens first, we will use the marginal $P(\X)$.
Therefore, the loss of using $P(\XY)$ includes the loss of using $P(\X)$.
Then, when \Y happens, since \X and \Y are independent, we will be using the marginal $P(\Y)$,
and the loss of using $P(\XY)$ also includes the loss of using $P(\Y)$.
Therefore, it appears that separately evaluating the marginal distributions $P(\X)$ and $P(\Y)$, should be the same as
evaluating $P(\XY)$ when \X and \Y are independent.
Hence, we may use
\[
    \D[\X] + \D[\Y].
\]
Separate evaluation of marginals is essential for the divergence measure that we will define shortly.

The KL divergence provides a method for measuring the agreement of a model about a specific distribution.
This agreement can be called convergence.

\begin{definition}[$\epsilon$-convergence]
    \label{def:e-convergence}
    In a probabilistic model \model, let \X and \E be discrete random variables.
    We say $\XIE$ (\X conditioned on \E) \emph{\econverges by $P \in \set{p}$},
    if for any probability measure $Q \in \set{p}$ we have:
    \begin{equation*}
        \Ex{Q}{\D[\XIe]} < \epsilon,
    \end{equation*}
    where $\Ex{Q}{\,\cdot\,}$ denotes expectation relative to $Q$ and $\D[\XIe]$ is the Kullback-Leibler divergence
    of $P(\XIe)$ from $Q(\XIe)$:
    \[
        \D[\XIe] = \sum_{\rval{x}} Q(\xIe) \ln \frac{Q(\xIe)}{P(\xIe)}\cdot
    \]
\end{definition}

$\epsilon$-convergence is closely related to the concept of Bayesian convergence, where the posterior converges to
the same distribution regardless of the prior distribution.

As an example, consider a scenario where we are trying to predict the outcome of a four-sided dice.
This dice is unfair, with each face marked by a binary number: $00$, $01$, $10$, and $11$.
Before making our prediction, we can roll the dice multiple times, with each roll being independent.

If we let \E be a vector representing the first 10 dice rolls and \X be the outcome of the next roll,
\XIE \econverges by $P_u$ for $\epsilon \geq 0.109$, where $P_u$ is the probability measure induced by considering a
uniform prior for the multinomial parameters.
If \E represents the first 20 outcomes, then we can choose any $\epsilon \geq 0.073$.

Here, the selected prior for the multinomial parameters affects the minimum value of $\epsilon$.
In other words, different probability measures in our model have different $\epsilon$-convergence properties.
The minimum $\epsilon$ gives us an assessment of the safety of a probability measure,
with respect to a particular \XIE conditional.
If this value is sufficiently small, any other probability measure in the model agrees that our distribution for \XIE is
safe, in the sense that it has a small expected divergence.

When we are interested in the distribution of several random variables, we may consider the sum of their KL
divergences, as we discussed earlier, to construct a measure that accounts for various random variables.

\begin{definition}[Query Set]
    A \emph{query set} is a set of pairs of discrete random variables, all defined on the same sample space.
\end{definition}

\begin{definition}[Divergence Measure]
    Let \Xpairs, be a query set in a probabilistic model \model.
    We define the \emph{divergence measure of $P \in \set{P}$ for \Qset} as follows:
    \begin{align}
        \label{eq:div_def}
        \Dx & = \supQ \DxExpand,                    \\
        \nonumber
        & = \supQ \sumixe Q(\assign[i]{x} , \assign[i]{e}) \ln
        \frac{Q(\assign[i]{x} \mid \assign[i]{e})}{P(\assign[i]{x} \mid \assign[i]{e})} \cdot
    \end{align}
\end{definition}

\begin{proposition}
    \XIE \econverges by $P$, if and only if \[\Dx[P][\{(\XE)\}] < \epsilon \cdot\]
\end{proposition}
Note that we usually simplify our notation and instead of writing the full assignment to a random variable, we only
write the assigned value.
For example, instead of $Q(\assign[i]{x} , \assign[i]{e})$ we write: $Q(\xe[i])$.

For a parametric model (as defined in Definition~\ref{def:param-model}), it seems reasonable to define a conditional
form of our divergence measure, conditioned on the value of the parameter.

To do so, in Equation~\ref{eq:div_def}, we replace normal expectation with conditional expectation.
There will be no longer a need to maximize over different distributions, since in a parametric model the parameter
fully determines the distribution of the sample space.

\begin{definition}[Conditional Divergence Measure]
    \label{def:cond_divergence}
    Assume \pmodel is a parametric model, and $\Xpairs$ is a query set.
    We define the \emph{conditional divergence measure of $P \in \set{P}$ for \Qset given $\paramValue$} as follows:
    \begin{align}
        \DxCond & = \DxCondExpand           \\
        & = \DxCondExpandFull \cdot
        \label{eq:cond_div_measure}
    \end{align}
\end{definition}

\begin{proposition}
    For the conditional divergence we have:
    \begin{itemize}
        \item $\begin{aligned}
                   \DxCond \geq 0
        \end{aligned}$
        \item $\begin{aligned}
                   \Dx = \supQ \intProd[q]{\left(\DxCond - \DxCond[Q] \right)}
        \end{aligned}$
    \end{itemize}
\end{proposition}

Subjectively, the probability measure that has the lowest divergence can be considered the ``safest'' choice
in a model.
It is safe in the sense that it has the best possible performance in the worst case, and more importantly, when
convergence is possible in a model, if we use this probability measure, we are guaranteed to
have a convergent distribution for our query set.

\begin{definition}[Central Distribution]
    In a model \model, we define the \emph{central distribution} for a query set \Qset, to be a probability
    measure $\central \in \set{P}$ which minimizes the divergence measure for \Qset:
    \begin{equation*}
        \central = \arg\minP \Dx\cdot
        \label{eq:center}
    \end{equation*}
\end{definition}

In this definition, the query set \Xpairs plays a crucial role.
Specifically, \Qset defines the set of random variables in which we are interested.
Each \X[i] represents a random variable whose distribution is important to us, and each \E[i] corresponds to the
evidence that is available for that variable.

However, if we choose our query set based on our specific experimental setup, it would violate the
principle of objectivity.
To avoid that, we have to use the same query set for any possible problem that could be
defined within an environment.

One natural way to do this is to let the query set include all potentially relevant conditional pairs.
For example, in our dice problem, we might let each \X[i] be the outcome of the $i$-th dice roll, and define each \E[i]
as a vector of the previous dice outcomes: $\E[i] = (\seqRv{i-1})$.
We may assume that our query set contains $n$ pairs.
The choice of $n$ should be based on our assumptions about the complexity of the problem.
It shows where we think the convergence becomes trivial and, if necessary, could approach infinity.

Now let us calculate the conditional divergence measure (Definition~\ref{def:cond_divergence}) for this query set:
\begin{align*}
    \DxCond & = \KlCond[][{\x[1]}][ ]                  \\
    & + \KlCond[][{\x[2]}][{\x[1]}]            \\
    & + \ \cdots                               \\
    & + \KlCond[][{\x[n]}][\seqRval{n-1}]\cdot
\end{align*}
After factoring out $\pi(\seqRval{n} \mid \paramValue)$ from every term:
\[
    \sum_{\seqRval{n}} \pi(\seqRval{n} \mid \paramValue) \ln
    \frac
    {\pi(\xIt[1]) \pi(\x[2] \mid \x[1], \paramValue)
    \ldots \pi(\x[n] \mid \seqRval{n-1}, \paramValue)}
    {P(\x[1]) P(\x[2] \mid \x[1]) \ldots P(\x[n] \mid \seqRval{n-1})},
\]
we use the chain rule for conditional probability:
\begin{equation}
    \DxCond = \KlCond[][\seqRval{n}][ ] \cdot
    \label{eq:joint_form}
\end{equation}

This means that we can alternatively define our query set as $\Qset = \{ (\XE) \}$,
where \X represents the vector of $n$ dice rolls: $\X = (\seqRv{n})$, and \E is a dummy certain event with $P(\E) = 1$.
This query set essentially corresponds to the entire sample set for $n$ dice rolls.
We usually simplify our notation and omit dummy events and write: $\Qset = \{\X\}$.

\newcommand{\XcondYseq}{\X[i] \mid \seqRv[Y]{i-1}}
\newcommand{\YcondYseq}{\Y[i] \mid \seqRv[Y]{i-1}}

The fact that our query set includes the entire sample set does not necessarily mean its central distribution has good
convergence properties for every random variable.
In our dice example, suppose we are interested in predicting the next dice outcome based on the number of previous
odd and even outcomes.
Let \Y[i] represent the least significant bit of the $i$-th dice roll, indicating whether the number is odd or even.
If we want our central distribution to exhibit good convergence properties for  $P(\XcondYseq)$,
we need to explicitly include these conditional pairs in our query set, in the same way we did for $\X[i]|\seqRv{i-1}$.

Intuitively, it seems that in any noninformative distribution we should have
\[
    P(\XcondYseq) = \frac{1}{2} P(\YcondYseq) \cdot
\]
For example, consider $P(\X[4] = 01 \mid 0,1,1)$, which represents the probability of observing $01$ on the 4th
roll, given that we have previously seen two odd numbers and one even number.
It is reasonable to expect that this probability equals $P(\X[4] = 11 \mid 0,1,1)$.
On the other hand,
\[
    P(\Y[4] = 1 \mid 0,1,1) = P(\X[4] = 01 \mid 0,1,1) + P(\X[4] = 11 \mid 0,1,1),
\]
which implies
\[
    P(\X[4] = 01 \mid 0,1,1) = P(\X[4] = 11 \mid 0,1,1) = \frac{1}{2} P(\Y[4] = 1 \mid 0,1,1) \cdot
\]
This suggests that adding $\YcondYseq$ conditionals to our query set,
would have the same effect as adding $\XcondYseq$ conditionals.

If we add conditionals of the form $\YcondYseq$ to our query set, we will get
$\Qset = \{\X, \Y\}$, where \X represents the sample space of $n$ rolls of a four sided dice,
and \Y represents the sample space of $n$ rolls of a two-sided odd-even dice or coin.

The central distribution for this new query set shows some interesting properties.
In this distribution, for smaller values of $i$, $P(\X[i] \mid \seqRv{i-1})$
tends to get somewhat combined with $P(\XcondYseq)$.
For example when $i = 2$, if the single observed dice roll is $01$, the posterior probabilities of
$00,10$ and $11$ are not equal, and $P(11 \mid 01) = 0.158$, which is slightly higher than
$P(00 \mid 01) = P(10 \mid 01) = 0.133$.
We can say that the central distribution identifies that $01$ and $11$ are both odd numbers.
